Some rules reject programs that could be accepted without loss of soundness, only so that a well-typed program is
also well typed under mypy (\secref{well-formedness}). \figref{mypy-restrictions} lists each restriction with mypy's
verdict and the rule that enforces it.

\begin{figure}
\centering
\footnotesize
\setlength{\tabcolsep}{4pt}
\begin{tabular}{>{\raggedright\arraybackslash}p{0.26\textwidth}>{\raggedright\arraybackslash}p{0.26\textwidth}>{\raggedright\arraybackslash}p{0.40\textwidth}}
\arrayrulecolor{gray!60}\hline
\rowcolor{grammarheadcolor}Construct & mypy & Mechanism \\
Call returning \pyinline{None} as an expression & Rejects any use of the value & \ruleName{call} requires a result other than \pyinline{None}; the call is a statement (\ruleName{call-stmt}) \\\rowgap
Body that may fall through at a return type such as \pyinline{int | None} & Reports a missing return & \ruleName{def-body} requires outcome $\Returns$ unless the return type is equivalent to \pyinline{None} \\\rowgap
\added{Assertion of an expression of type \pyinline{Literal[False]} other than the literal} & \added{Reports a missing return} & \added{\ruleName{assert-false} applies to the literal \pyinline{False} only} \\\rowgap
Bare \pyinline{return} at such a return type & Reports a missing return value & \ruleName{return-none} requires the return type equivalent to \pyinline{None} \\\rowgap
Match whose cases cover a class field by field & Narrows the scrutinee by whole classes, so reports a missing return & Exhaustiveness decided by the remaining type (\ruleName{match}) \\\rowgap
Pattern variable after a case excluding a field value & Binds at the field type & Bindings taken from the remaining type (\ruleName{match-case}, \ruleName{check-var}) \\\rowgap
\pyinline{list[bool]} below \pyinline{list[int]} & Rejects; list and dictionary types invariant & \ruleName{subty-list} and \ruleName{subty-dict} require equivalent components \\\rowgap
Declaration of a name already a parameter, imported name, pattern variable or declaration of the scope & Rejects the redeclaration & Context $\scopeS(Y, s)$ of the body undefined unless the names are pairwise distinct
\end{tabular}
\caption{Restrictions imposed for compatibility with mypy}
\label{fig:mypy-restrictions}
\end{figure}


\begin{AddedBlock}
mypy may also reject a well-formed program because of a known bug in mypy. There is one such bug,
python/mypy\#22202.\footnote{\url{https://github.com/python/mypy/issues/22202}} When deciding whether a match is
exhaustive, mypy does not use a Boolean value known from literals in the scrutinee, such as \pyinline{False} or
\pyinline{(False,)[0]}, although it uses an integer or string value known in the same way. It therefore reports a
missing return for the following well-formed function:

\begin{python}
def f() -> int:
    match False:
        case False:
            return 0
\end{python}

\noindent The conformance suite keeps such programs apart from the other semantically valid tests, and checks that
mypy still rejects them.
\end{AddedBlock}
